% !TEX root = report_template.tex
\documentclass{article}
\setlength{\parskip}{1ex}
\title{ DESeq2 Report: {{comparison_clean}} }
\author{ICBI}
\usepackage{graphicx}
\usepackage{longtable}
\usepackage{geometry}
\usepackage{helvet}
\usepackage{natbib}
\usepackage{placeins}
\usepackage{booktabs}
\usepackage{amsmath} % for better handling of underscores in words
\renewcommand{\familydefault}{\sfdefault}
\geometry{margin=1in}

\begin{document}

\maketitle

\section*{Introduction}
This report summarizes the results of the comprehensive analysis of RNA-seq data,
encompassing differential expression analysis, quality control visualization, and functional
enrichment analysis.

\section*{Methods}
\begin{longtable}{ll}
    \hline
    Category & Details \\
    \hline
    Comparison & {{comparison_clean}} \\
    Paired Analysis & {{#paired}} {{paired}} {{/paired}} {{^paired}} No {{/paired}} \\
    Covariates & {{#covariates}} {{covariates}} {{/covariates}} {{^covariates}} None {{/covariates}} \\
    Batch Effects Considered & {{batchEffects}} \\
    Batch & {{batch_col}} \\
    \hline
\end{longtable}

\section*{Result File Descriptions}
\begin{itemize}
    \item \textbf{`raw\_counts.tsv`}: Contains the raw gene count table across all samples.
    \item \textbf{`all\_samples\_PCA.png`}: A PCA plot displaying all samples to provide an overview of sample distribution.
    \item \textit{multiqc\_report.html}: HTML report from MultiQC summarizing data quality across samples.
    \item \textbf{Feature and Gene Detection Files}:
    \begin{itemize}
        {{#has_deg}}
        \item \textbf{`\textasteriskcentered{}\_biotype\_counts.png` \& `\_biotype\_counts.tsv`}: Lists and visualizations of the most differentially expressed genomic features by type (e.g., protein-coding).
        {{/has_deg}}
        \item \textbf{`\textasteriskcentered{}\_number\_of\_detected\_genes.png (tsv)`}: Bar plot of the number of detected genes (min 10 reads) per sample.
        \item \textbf{`\textasteriskcentered{}\_detectedGenesNormalizedCounts\_min\_10\_reads\_in\_one\_condition.tsv`}: Contains normalized counts of genes detected with at least 10 reads in one condition.
        \item \textbf{`\textasteriskcentered{}\_detectedGenesRawCounts\_min\_10\_reads\_in\_one\_condition.tsv`}: Contains raw counts of genes detected with at least 10 reads in one condition.
    \end{itemize}
    {{#has_deg}}
    \item \textbf{Differential Expression Analysis}:
    \begin{itemize}
        \item \textbf{`\textasteriskcentered{}\_IHWsigFCgenes\_2\_fold.xlsx (tsv)`}: Significant differentially expressed genes with fold change \textgreater 2 and FDR \textless 0.1.
        \item \textbf{`\textasteriskcentered{}\_IHWsigGenes.xlsx (tsv)`}: Significant differentially expressed genes with FDR \textless 0.1 (no fold change cutoff).
        \item \textbf{`\textasteriskcentered{}\_IHWallGenes.xlsx (tsv)`}: All genes in each comparison with no filters applied.
    \end{itemize}
    {{/has_deg}}
    \item \textbf{Visualization Files}:
    \begin{itemize}
       {{#has_deg}}
        \item \textbf{`\textasteriskcentered{}\_volcano\_padj.png`}: Volcano plot showing gene expression differences between conditions, with selected genes highlighted.
        {{/has_deg}}
        \item \textbf{`\textasteriskcentered{}\_PCA.png`}: PCA plot to visualize clustering by condition.
    \end{itemize}
    {{#has_deg}}
    \item \textbf{Gene Set Enrichment Analysis (GSEA) and Overrepresentation Analysis (ORA)}:
    \begin{itemize}
        \item \textbf{GSEA and ORA Dot Plots}: Visualization files and analysis results for gene-set enrichment and overrepresentation analysis, including GO terms (BP, MF, CC), KEGG pathways, Reactome pathways, and WikiPathways.
        \item \textbf{GSEA*.tsv and ORA*.tsv}: Table files containing results for each gene set.
        \item \textbf{Top GO Analysis}:
        \textbf{topGO\_IHWsig*.xlsx (tsv)}: Lists the statistically overrepresented GO terms in the categories Biological Process (BP), Molecular Function (MF), Cellular Component (CC).
    \end{itemize}
    {{/has_deg}}

\end{itemize}

\pagebreak
\section*{Differential Gene Expression with DESeq2}

DESeq2 is an R package for identifying differentially expressed genes from RNA-seq data.
It models count data with the negative binomial distribution, normalizes for library size differences,
and uses generalized linear models to test for significant changes between groups.

Coupled with Independent Hypothesis Weighting (IHW), DESeq2 improves statistical power while controlling
the false discovery rate by weighting genes based on independent covariates like their average expression level.
This combination provides a robust and widely-used method for differential gene expression analysis.

The output of DESeq2 includes log2 fold changes, shrunk log2 fold changes, p-values, and adjusted p-values for each gene, allowing
to identify and prioritize genes with significant expression differences between conditions.

\textbf{Important Considerations on Log2 fold change shrinkage}:

Log2 fold change shrinkage is a post-processing step in differential expression analysis that improves the accuracy and interpretability
of results by borrowing information across genes to refine individual estimates. It's particularly important for genes with low counts
or high variability, leading to more reliable visualizations and gene rankings.

\begin{itemize}
  \item Noise in Lowly Expressed Genes: Genes with low read counts are more susceptible to noise and variability. This can lead to inflated and unreliable log2 fold change estimates.
  \item Shrunk log2 fold changes (log2FoldChange\_shrink) are generally preferred for visualization and ranking genes, as they are more stable, especially for genes with low counts.
  \item The un-shrunk fold change is still important because it is used for the test statistic calculation.
  \item The adjusted p-value (padj) is crucial for determining which genes are considered significantly differentially expressed.
  \item log2FoldChange\_shrink is not used for the test statistic.
  \item Impact on Interpretation: Inflated log2 fold changes for lowly expressed genes can be misleading, making them appear more important than they actually are. This affects
  downstream analyses like ranking genes for further investigation or creating visualizations like MA-plots and volcano plots.
\end{itemize}

{{#has_deg}}

\begin{table}[h!]
  \centering
  \small % Reduce font size for the entire table
  \caption{Description of columns in a DESeq2 output file: {{comparison_clean}} }
  \label{tab:deseq2_output}
  \begin{tabular}{|p{0.25\linewidth}|p{0.68\linewidth}|}
  \toprule
  \textbf{Column Name} & \textbf{Description} \\
  \midrule
  \textbf{gene\_id} & The unique identifier for the gene (Ensembl ID) \\
  \textbf{baseMean} & The average normalized read count across all samples, after normalization by library size factors. \\
  \textbf{log2FoldChange} & The estimated log2 fold change between the two conditions being compared. A positive value indicates higher expression in the first condition. \textbf{This is the unshrunk version.} \\
  \textbf{lfcSE} & The standard error associated with the \textbf{unshrunk} \texttt{log2FoldChange} estimate. \\
  \textbf{stat} & The Wald statistic, calculated as \texttt{log2FoldChange} / \texttt{lfcSE}. \\
  \textbf{pvalue} & The raw p-value resulting from the Wald test, representing the statistical significance of the observed fold change. \\
  \textbf{padj} & The p-value adjusted for multiple testing (IHW), controlling the false discovery rate (FDR). \\
  \textbf{log2FoldChange\_shrink} & The \textbf{shrunk} log2 fold change. This value is shrunk towards zero from the original \texttt{log2FoldChange}, especially for genes with low counts or high dispersion. \texttt{normal} or \texttt{apeglm} shrinkage can be applied. \\
  \textbf{lfcSE\_shrink} & The standard error associated with the \textbf{shrunk} \texttt{log2FoldChange\_shrink} estimate. \\
  \textbf{gene\_name} & The official gene symbol \\
  \textbf{genes\_description} & A description of the gene function, if available and added from an annotation database. \\
  \textbf{gene\_type} & Biotype \\
  \bottomrule
  \end{tabular}
\end{table}

{{/has_deg}}


\FloatBarrier

\subsection*{DGE plots}
{{#pca_plots}}
    {{#.}}
    \begin{figure}[htbp]
        \centering
        \begin{minipage}[t]{0.48\textwidth}
            \centering
            {\includegraphics[width=\textwidth]{ {{image_path}} }}
            \caption{ {{caption}} }
        \end{minipage}
        {{#has_pair}}
        \hspace{0.01\textwidth}
        \begin{minipage}[t]{0.48\textwidth}
            \centering
            {\includegraphics[width=\textwidth]{ {{pair_path}} }}
            \caption{ {{pair_caption}} }
        \end{minipage}
        {{/has_pair}}
    \end{figure}
    {{/.}}
{{/pca_plots}}

{{#no_deg}}
    \textbf NO differntially expressed genes found when comparing {{comparison_clean}} 
{{/no_deg}}

{{#has_deg}}

{{#expr_plots}}
    {{#.}}
    \begin{figure}[htbp]
        \centering
        \begin{minipage}[t]{0.48\textwidth}
            \centering
            {\includegraphics[width=\textwidth]{ {{image_path}} }}
            \caption{ {{caption}} }
        \end{minipage}
        {{#has_pair}}
        \hspace{0.01\textwidth}
        \begin{minipage}[t]{0.48\textwidth}
            \centering
            {\includegraphics[width=\textwidth]{ {{pair_path}} }}
            \caption{ {{pair_caption}} }
        \end{minipage}
        {{/has_pair}}
    \end{figure}
    {{/.}}
{{/expr_plots}}

\FloatBarrier
\pagebreak
\section*{Functional Enrichment Analysis: ORA \/ GSEA}

The over-representation test (ORA) takes the list of deferentially expressed genes (as defined by the FDR- and fold-change cutoff)
and checks if it contains more genes from a certain gene-set (e.g. GO-term, KEGG pathway) than would be expected
by random chance.

Gene-set-enrichment-analysis (GSEA) ranks genes from most upregulated to most downregulated, and checks if genes from a certain gene-set
tend to be more at the top, or more at the bottom of this list.

The ORA test results are straightforward to interpret and, in most cases, the preferred flavor of gene-set analysis.
GSEA is superior when very few genes are differentially expressed, as it also takes into account more subtle,
but coordinated gene expression changes that are statistically significant only when regarding the gene set as a whole,
but not on the level of single genes.


\subsection*{Over-Representation Analysis (ORA)}

ORA was performed using clusterProfiler and topGO packages to identify over-represented pathways and Gene Ontology (GO) terms
among the significantly differentially expressed genes.

\textbf{Gene Set and Pathway tested for enrichment:}

\begin{itemize}
  \item KEGG pathways
  \item Reactome pathways
  \item WikiPathways
  \item GO terms:
  \begin{itemize}
    \item BP: Biological Process
    \item MF: Molecular Function
    \item CC: Cellular Component
  \end{itemize}
\end{itemize}

\FloatBarrier
\subsection*{ORA plots of the top 30 significantly overrepresented gene sets}
{{#figures_ORA}}
    {{#path}}
    \begin{figure}[htbp]
        \centering
        \hspace*{-3cm}
        {\includegraphics[width=\textwidth]{ {{path}} }}
        \caption{ {{caption}} }
    \end{figure}
    {{/path}}
{{/figures_ORA}}

\FloatBarrier
\subsection*{Gene Set Enrichment Analysis (GSEA)}

GSEA was performed using clusterProfiler to assess whether pre-defined gene sets (e.g., pathways) are enriched at the top or bottom
of a ranked list of genes, where the ranking is based on the differential expression test statistic.

\textbf{Gene Set and Pathway tested for enrichment:}

\begin{itemize}
  \item KEGG pathways
  \item Reactome pathways
  \item WikiPathways
  \item GO terms:
  \begin{itemize}
    \item BP: Biological Process
    \item MF: Molecular Function
    \item CC: Cellular Component
  \end{itemize}
\end{itemize}

\FloatBarrier
\subsection*{GSEA plots of the top 30 significantly enriched gene sets}
{{#figures_GSEA}}
    {{#path}}
    \begin{figure}[htbp]
        \centering
        \hspace*{-3cm}
        {\includegraphics[width=\textwidth]{ {{path}} }}
        \caption{ {{caption}} }
    \end{figure}
    {{/path}}
{{/figures_GSEA}}

{{/has_deg}}


\FloatBarrier
\nocite{*}
\bibliographystyle{plainnat}
\bibliography{ {{bib_file}} }

\end{document}
